\section{Conclusion}\label{sec:conclusion}

xlog is a CUDA-native logic programming engine with one typed language and
provider-owned runtime services for deterministic, probabilistic, epistemic,
and neural-symbolic use. A neural network is an ordinary predicate; exact query
gradients flow through the probabilistic knowledge-compilation route, while the
other routes reuse typed device-buffer interfaces without claiming identical
differentiability. The final smoothed probabilistic circuit is certified against
its back-end source formula before it is cached or evaluated. Solver proof traces
are checked on the GPU, and the host observes bounded status and error scalars.

Residency claims are likewise route-specific. Ordinary relational and exact
execution are host-orchestrated over device data planes. Eligible recursion and
resident Monte Carlo use certified cores that record zero tracked transfers and
then return one authoritative terminal receipt. Byte exhaustion, fixed-capacity
exhaustion, and nonconvergence remain separate typed outcomes.

The measurements describe a system that pays off where the symbolic workload
is large and repeated: reusing a verified circuit across training iterations,
and avoiding the intermediate blowup of binary-join plans on skewed cyclic
queries. Against external engines the picture is mixed: clear
gains in bounded-memory join evaluation and in per-epoch neurosymbolic
training, and a clear loss to a mature CPU compiler on small exact-inference
programs.

What they do not give is a general ranking of xlog against other systems. Each
comparison covers one engine per reasoning mode on one class of program, the
rule-induction results rest on two public corpora, and the epistemic and
rule-learning surfaces are young. The evidence supports the architecture, not
a claim of universal superiority.

The direction is the one the architecture already implies: out-of-core and
multi-GPU execution to lift the memory ceiling, deeper residency for the
epistemic routes, and structure learning that runs on the stream rather than
over it.
